\documentclass[twoside,openright]{xduugthesis}

\usepackage{float}

\xdusetup{
  style / font-type = font,
  style / latin-font = tacn,
  style / cjk-font = fandol,
  info / title = {面向标注稀缺的医学图像自适应分割\\理论与方法研究},
  info / author = {杨尚霖},
  info / department = {计算机科学与技术学院},
  info / major = {计算机科学与技术},
  info / supervisor = {权义宁},
  info / class-id = {2203018},
  info / student-id = {22009290006},
  info / abstract = {abstract-zh.tex},
  info / abstract* = {abstract-en.tex},
  info / keywords = {医学图像分割，半监督学习，标注稀缺样本，自适应分割},
  info / keywords* = {Medical Image Segmentation, Semi-supervised Learning, Annotation-scarce Samples, Adaptive Segmentatio},
  info / bib-resource = {references.bib}
}
\usepackage{amsmath}   % 解决 \operatorname, \exp 等报错
\usepackage{amssymb}   % 解决 \mathbb{R}, \mathbb{N} 等报错
\usepackage{booktabs}  % 解决 \toprule, \midrule 等表格报错

\begin{document}


\chapter{绪论}

\section{国内外研究现状}
医学图像分割是计算机辅助诊断系统中的关键技术之一，近年来随着深度学习技术的快速发展，基于全监督学习的分割方法在多个公开数据集上取得了显著成果。然而，这些方法高度依赖大量精确标注的数据，而在医学影像领域，高质量的标注数据通常稀缺且获取成本高昂，严重制约了模型在真实临床场景中的推广与应用。同时由于分割任务的不同，往往需要采用不同的特定分割网络，这大大增加了分析成本以及训练成本，需要一种能融合多尺度信息的通用网络框架。

\subsection{国内研究现状}
国内研究团队在医学图像分割领域已取得一系列进展。例如：DTC\cite{luo2021semi}采用多个独立的分割网络构成协同训练小组，让它们互为师生、通过视觉摘要和预测分歧相互监督，从而高效挖掘无标签数据中的互补信息；SASSNet\cite{li2020shape}将形状先验知识显式地建模为一个并行的形状流网络，通过强制分割预测与学习到的形状表示在潜在空间中对齐，确保模型对无标签数据的输出不仅在像素级别准确，同时符合解剖结构的整体合理性。

CA-Net\cite{gu2020ca}通过空间、通道及尺度三重注意力机制逐步细化分割结果，并结合课程学习策略从易到难学习病灶特征，显著提升了胰腺、肝脏等复杂解剖结构的分割稳定性；TransFusion\cite{bai2022transfusion}通过交叉注意力机制动态整合不同成像模态的特征表示，在病灶边界模糊、模态信息不匹配情况下仍能保持分割一致性，显著提升了多模态数据协同分析的鲁棒性。

\subsection{国外研究现状}
国外在医学图像分割领域的研究起步较早，尤其在自适应学习、半监督与自监督学习方面进展显著。例如：Mean Teacher\cite{tarvainen2017mean}通过平均学生模型的权重来生成更稳定的教师模型，利用一致性损失提升半监督学习的效果；MixMatch\cite{berthelot2019mixmatch}将多张有标签和无标签数据在图像和标签空间同时进行混合，并对无标签数据产生强化后的“猜测标签”，从而在一个损失中统一优化所有数据。

UNETR\cite{hatamizadeh2022unetr}首次将标准的Vision Transformer直接扩展为3D医学图像分割的编码器骨干网络，通过将体积图像分割为三维补丁序列并输入Transformer，以直接建模图像中任意两个体素之间的长距离全局依赖，再通过一个专用的卷积解码器逐步上采样和恢复空间细节，实现了对复杂解剖结构的全局上下文与局部空间一致性的统一建模；nnUNet\cite{isensee2018nnu}通过一个高度自动化的框架，系统化地优化数据预处理、训练策略和后处理流程，使标准的U-Net架构在不同医学图像数据集上都能达到或超越许多定制化模型的性能，强调了工程化流程的重要性。

\section{背景介绍}
随着深度学习理论与技术的飞速发展，医学图像分割作为计算机辅助诊断的核心环节之一，对于病灶定位、器官重建以及手术规划具有至关重要的意义。然而，当前的医学图像分割研究面临着两大核心挑战，这也是本文致力于解决的关键问题。

首先，高质量的医学图像数据标注成本极高，需要依赖资深医生的专业知识进行手工勾画，耗时耗力。这导致了标注数据稀缺，标注代价昂贵等问题，严重限制了深度学习模型的泛化能力。如何利用大量未标注的医学图像数据，通过半监督学习策略来挖掘数据本身的潜在信息，成为了当前研究的热点与难点。传统的半监督方法往往忽视了模型在预测过程中的可靠性评估，导致无标签数据中的噪声干扰了模型训练。

其次，现有的主流分割网络（如U-Net）主要基于卷积神经网络（CNN），其局部感受野特性限制了对长距离依赖关系的建模能力，难以根据所接受的病理特征进行自适应推理，并且在处理具有复杂背景和模糊边界的医学图像时，容易出现细节丢失或误分割。近年来，Transformer架构在自然语言处理领域取得巨大成功后，逐渐被引入计算机视觉领域，其强大的全局建模能力为解决上述问题提供了新的思路，但如何将其有效融入医学图像分割框架，并与传统的卷积网络优势互补，仍需深入探索。

基于上述背景，本文聚焦于医学图像分割任务，从两个关键维度展开研究：一是引入认知不确定性理论，通过改进Mean Teacher等半监督框架，使模型能够识别并抑制无标签数据中的不可靠预测，从而提升模型在数据稀缺情况下的鲁棒性；二是探索Transformer架构的引入机制，构建能够同时捕获局部细节与全局上下文信息的混合模型，以增强分割边界精度和对复杂结构的解析能力。通过这两条技术路线的并行研究，旨在为高精度、低标注成本的医学图像分析提供有效的解决方案。

接下来我将介绍几个基础概念，以进行对本文内容的补充与完善。

\subsection{机器学习范式：监督、无监督与半监督}
在机器学习领域，根据数据属性和学习目标的不同，可以将学习方法主要划分为三大基本范式：监督学习、无监督学习和半监督学习。这三种范式构成了数据挖掘与模式识别的理论基石，各自适用于不同的应用场景，并解决不同类型的现实问题。

监督学习是机器学习中最成熟、应用最广泛的范式。其核心思想是利用带有标签的训练数据集来学习一个从输入到输出的映射函数。所谓“标签”，即数据的已知结果或类别（如价格预测任务中某个房子具体的价格，或图像分类任务中的“猫”或“狗”）。在训练过程中，模型通过不断比较预测输出与真实标签之间的差异，调整内部参数以最小化误差。而在本文所进行的研究中，监督学习主要是作为性能基准，用于评估和对比半监督学习在利用无标签数据后所带来的性能提升效果。

与监督学习相对，无监督学习处理的是没有标签的原始数据。其主要目标不是进行预测，而是探索数据内部隐藏的结构、模式或分布。无监督学习不需要预先定义输出结果，而是通过算法自动发现数据的内在关联。在本文中由于没有对无监督学习的应用，便不多进行赘述。

半监督学习则介于上述两者之间，它试图结合少量的标注数据和大量的未标注数据进行模型训练。在现实世界中，获取大量精确标注的数据往往成本高昂且耗时（如医学影像标注），而未标注的数据则相对容易获得。半监督学习通过利用未标注数据的分布信息，辅助标注数据构建更鲁棒、泛化能力更强的模型。其基本假设通常是：数据具有一定的内在结构（如聚类假设或流形假设），即相近的数据点倾向于拥有相同的标签，以此对模型进行训练。本文在第三章中将对该方法进行具体的详细阐述与实验验证。”

\subsection{深度学习与神经网络概述}
深度学习作为机器学习的一个重要分支，其核心机制是利用深层神经网络来模拟人脑处理信息，从而实现对数据的多层次抽象与特征学习。它在传统神经网络的基础上进行了继承与发展：神经网络构成了深度学习的基础架构，而深度学习则通过引入更深的网络层次和更复杂的非线性变换，赋予了模型提取更高层次、更复杂特征的能力。

人工神经网络（Artificial Neural Network, ANN）受生物神经系统启发，主要由大量相互连接的“神经元”节点组成。一个基本的神经元模型接收输入特征向量x，通过加权求和并加上偏置项b，再经过非线性激活函数产生输出：
\begin{equation}
    y = f\left( \sum_{i=1}^{n} w_i x_i + b \right)
\end{equation}

其中$w_i$为连接权重，代表输入特征的重要性。通过调整权重和偏置，神经网络能够学习输入与输出之间的复杂映射关系。典型的网络结构包含输入层、隐藏层和输出层；随着隐藏层数量的增加，网络的表达能力显著增强。

深度学习模型的训练是一个通过数据驱动来自动调整内部参数的迭代优化过程，其核心目标是最小化模型预测输出与真实标签之间的误差。该过程主要包含以下四个关键步骤：

前向传播：数据从输入层进入网络，逐层经过加权求和与非线性激活函数的计算，最终在网络的输出层产生预测结果。这一过程将输入数据映射为模型的预测输出。

损失计算：利用损失函数来量化模型预测值与真实标签之间的差距。损失函数的值被称为“损失”（Loss）或“误差”（Error），训练的目标即是最小化这一数值。

反向传播：这是训练深层网络的核心步骤。它从输出层开始，将损失值沿着网络结构反向传递。通过计算损失函数对每个网络参数（权重w和偏置b）的偏导数，从而明确参数更新的方向。

参数更新：在获取了参数的梯度后，优化器会使用这些信息来更新网络参数。最基础的更新规则是梯度下降法，其核心思想是让参数沿着梯度的反方向（即损失减小最快的方向）进行调整。更新公式可表示为：
\begin{equation}
    w_{\text{new}} = w_{\text{old}} - \eta \cdot \nabla_w
\end{equation}
​
其中，$w_{\text{new}}$ 和 $w_{\text{old}}$ 分别代表更新后和更新前的权重，$\nabla_w$ 是损失函数对权重的梯度，而 $\eta$ 则是学习率。学习率是一个关键的超参数，它控制着参数更新的步长。

通过反复执行上述“前向传播 → 损失计算 → 反向传播 → 参数更新”的循环，模型能够逐步学习到数据中的内在规律，使其预测能力不断提升。

\section{论文主要工作}
在半监督医学图像分割领域，如何有效利用大量未标注数据以缓解对昂贵标注数据的依赖，已成为当前研究的热点与难点。深入挖掘未标注数据中的隐含信息有利于提升模型在复杂场景下的泛化能力。然而，现有的半监督方法在训练过程中往往忽略了模型对预测结果的可信度评估，导致在迭代过程中容易产生错误累积（即确认偏差）或对噪声过度信赖。同时，传统的卷积神经网络（CNN）在捕捉长距离依赖关系方面存在局限，而直接引入Transformer架构虽然增强了全局建模能力，但如何高效地融合局部细节与全局语义信息仍是一个挑战。

对于上述半监督分割算法中的问题，我们需要对模型的不确定性进行建模以指导训练过程，并对不同的网络架构进行消融实验以验证其有效性，所以本文的主要工作是： 

（1）引入认知不确定性的半监督学习范式：将认知不确定性引入Mean Teacher模型中，利用蒙特卡洛Dropout估计预测的不确定度，并将其作为权重融入一致性损失函数，从而抑制高不确定性样本对模型更新的负面影响

（2）构建并对比多种CNN-Transformer混合架构：在U-Net的基础上结合Transformer，实现了Swin-Unet和TransUnet两种不同的结合方式，并在此基础上进一步引入DA-TransUnet的跳跃连接注意力机制，对这四种模型（U-Net baseline, Swin-Unet, TransUnet, DA-TransUnet）进行了系统的对比分析；

（3）定量评价模型性能：在标准测试集上对算法结果进行定量评价（如Dice系数、Hausdorff距离）并对不同的模型结构进行评价。

\section{论文章节安排}
本论文围绕半监督医学图像分割中的不确定性建模与网络架构优化展开研究，旨在提升模型在有限标注数据下的鲁棒性与分割精度。全文共分为五个章节，具体的组织结构安排如下：

第一章：绪论。本章首先介绍了机器学习的基本范式，重点阐述了半监督学习在医学图像分析领域的研究背景与应用价值。随后概述了深度学习与神经网络的基本概念，明确了本文的主要研究工作，并给出了论文的整体章节安排。

第二章：相关理论基础。本章系统地介绍了支撑本文研究工作的核心理论。首先详细阐述了认知不确定性理论，包括其定义、成因以及常用的度量方法；接着深入剖析了Transformer架构的原理，重点解析了自注意力机制，并介绍了Vision Transformer与Swin Transformer在视觉任务中的应用，为后续模型设计奠定理论基础。

第三章：基于认知不确定性的半监督方法研究。本章聚焦于半监督学习框架的改进。首先介绍了Mean Teacher框架的基本原理，随后重点阐述了UA-MT（不确定性感知的Mean Teacher）框架的构建过程及不确定性引入机制。通过设计相应的实验，利用特定数据集与评价指标对算法性能进行了对比验证与结果分析。

第四章：基于Transformer的多尺度自适应表示学习模型。本章详细介绍了网络模型的搭建与融合策略。首先确定了U-Net作为基础模型，随后探讨了Transformer的引入方式，具体涵盖了Swin-Unet、TransUnet以及改进的DA-TransUnet模型。最后通过系统的实验设计，对不同模型在相同环境下的表现进行了评估与对比。

第五章：总结与展望。本章对全文的研究工作进行了全面总结，归纳了论文所取得的主要成果。同时，客观指出了当前研究中存在的不足之处，并对未来可能的研究方向及应用前景进行了探讨。

\chapter{相关理论基础}
为有效解决绪论中提出的模型预测可信度评估问题以及模型在长距离依赖捕捉与任务自适应学习方面不足的问题，本章系统性地阐述了研究所依赖的核心理论框架，主要是相关的理论基础。一方面，通过梳理认知不确定性理论的定义、成因及度量手段，确立了评估模型预测结果置信度的理论依据；另一方面，通过对Transformer架构的理解与剖析，也能更深刻的理解Transformer能解决相关问题的内在原因。上述理论基础的建立，不仅完善了本文的知识体系，也为后续章节的创新性研究提供了关键的理论支撑。

\section{认知不确定性理论}
在深度学习与统计领域，模型的预测不仅需要准确性，更需要具备对自身预测能力的“自知之明”，即能认识到自己有多大的把握确认自己的判断。认知不确定性正是用来刻画这种“无知程度”的概念，它源于模型对输入数据知识的缺乏，反映了参数层面的不确定性。本节将系统阐述其定义与主流度量方式。

\subsection{认知不确定性的定义与成因}
认知不确定性，又称知识不确定性，是指由于训练数据覆盖不足、模型表达能力受限等原因导致的，会使模型在面对未知分布样本时产生的预测歧义性。其本质在于模型参数所具有的不确定性——即在给定数据下，存在多个参数配置均能较好拟合观测，但对新样本给出不同预测。从贝叶斯视角看，若将模型参数w视为随机变量，认知不确定性即来源于后验分布 $p(w\mid D)$的弥散程度：分布越分散，说明数据对参数的约束越弱，模型认知越模糊。

其成因主要包括三方面：一是数据稀缺性，在输入空间的某些区域缺乏足够标注样本，导致模型无法学习到对应有效的决策边界；二是模型表达局限性，即使数据充足，固定结构的模型（如结构固定的网络模型）也可能无法捕捉复杂的非线性映射；三是路径依赖，随机梯度下降等算法可能收敛到不同的局部最优解，加剧预测的不一致性。。

\subsection{认知不确定性的度量方法}
对认知不确定性的量化主要依赖于对模型参数分布的近似推断。主流方法是通过采样来估计预测分布的离散程度。其中，蒙特卡洛Dropout作为一种高效实用的方法，通过在测试阶段多次启用Dropout进行前向传播，生成一组 Softmax 概率向量 $\{\mathbf{p}_t\}_{t=1}^T$，进而估计不确定性：
$$
\mu_c = \frac{1}{T} \sum_{t=1}^T \mathbf{p}_t^c, \quad u = -\sum_{c} \mu_c \log \mu_c
$$

该方法的理论基础在于将Dropout视为对贝叶斯后验的变分近似，从而将认知不确定性与随机不确定性解耦。
另一种重要范式是深度集成，通过训练多个独立初始化的模型并集成其预测，由于本文并没有使用该方法，故不进行详细介绍。

\section{Transformer架构原理}

\subsection{自注意力机制}

自注意力机制（Self-Attention Mechanism）是 Transformer 架构的核心，它通过计算序列中不同位置元素之间的关联权重，实现对全局依赖关系的建模。其本质是将输入序列映射为查询（Query，$Q$）、键（Key，$K$）和值（Value，$V$）三个矩阵，通过点积运算衡量查询与键的相似性，再利用 Softmax 函数归一化得到注意力权重，最终加权求和值矩阵以生成输出。

为缓解点积结果过大导致的梯度消失问题，通常引入缩放因子 $d_k$（$d_k$ 为键向量的维度），其核心计算公式为：
\begin{equation}
\operatorname{Attention}(Q, K, V) = \operatorname{Softmax}\left(\frac{QK^{T}}{\sqrt{d_{k}}}\right) V
\end{equation}

为进一步提升模型的表征能力，Transformer 采用多头注意力机制，将 $Q$、$K$、$V$ 分别投影到 $h$ 个不同的子空间（“头”）中并行计算注意力，再将各头的输出拼接后线性变换。公式化表达为：
\begin{equation}
\operatorname{MultiHead}(Q, K, V) = \operatorname{Concat}(\operatorname{head}_{1}, \operatorname{head}_{2}, \dots, \operatorname{head}_{h}) W^{O}
\end{equation}

其中，第 $i$ 个头的计算为：
\begin{equation}
\operatorname{head}_{i} = \operatorname{Attention}(QW^{Q}_{i}, KW^{K}_{i}, VW^{V}_{i})
\end{equation}

$W^{Q}_{i}$、$W^{K}_{i}$、$W^{V}_{i}$ 为投影矩阵，$W^{O}$ 为输出投影矩阵。

\subsection{Vision Transformer与Swin Transformer}
Vision Transformer（ViT）将 Transformer 架构直接迁移至图像领域，核心思想是将图像分割为固定大小的图像块（Patch），并将每个图像块展平后线性投影为“视觉词元”（Visual Token），再添加位置编码（Positional Encoding）以保留空间信息。这些词元序列被输入标准 Transformer 编码器，通过自注意力机制捕捉全局特征。

ViT 的输入处理可形式化为：设输入图像尺寸为 $H \times W \times C$（高、宽、通道数），图像块尺寸为 $P \times P$，则图像块数量为 $N = \frac{HW}{P^2}$，每个图像块的展平向量为 $x_p \in \mathbb{R}^{P^2 \cdot C}$，经线性投影后得到词元：
\begin{equation}
    x_i = W_p x_p + E_{\text{pos},i}
\end{equation}

其中 $W_p$ 为投影矩阵，$E_{\text{pos},i}$ 为位置编码。

Swin Transformer 针对 ViT 计算复杂度高的问题，提出滑动窗口机制，将自注意力计算限制在局部窗口内，同时通过窗口滑动实现跨窗口的信息交互。其核心是窗口注意力（Window Attention）和移位窗口（Shifted Window）：

\begin{itemize}
    \item \textbf{窗口注意力}：将特征图划分为不重叠的局部窗口（如 $M \times M$），在每个窗口内计算自注意力，计算复杂度由 $O(N^2)$ 降至 $O(M^2 \cdot \frac{N}{M^2}) = O(N)$（$N$ 为词元数量，$M$ 为窗口尺寸）。
    \item \textbf{移位窗口}：在相邻 Transformer 层中，将窗口沿空间维度滑动 $M/2$ 个像素，使不同窗口的词元能够交互，避免局部窗口的孤立性。
\end{itemize}

Swin Transformer 还引入层次化特征图，通过逐步合并图像块构建多尺度特征，适配图像分类、检测、分割等下游任务。其窗口注意力的计算公式为：
\begin{equation}
    \operatorname{WindowAttention}(Q, K, V) = \operatorname{Softmax}\left(\frac{QK^T}{\sqrt{d_k}} + B\right) V
\end{equation}

其中，$B$ 为相对位置编码（Relative Position Bias），用于建模窗口内词元的相对空间位置关系。

\chapter{基于认知不确定性的半监督方法研究}
\section{Mean Teacher框架}
\subsection{Mean Teacher基本原理}
Mean Teacher框架是一种基于一致性正则化的半监督学习方法，其核心思想在于通过指数移动平均的方式更新教师模型的权重，并利用学生模型与教师模型之间的协同训练机制来提升模型性能。

\begin{figure}
    \centering
    \includegraphics[width=1.0\linewidth]{project_30_e2d9f79f.png}
    \caption{Mean Teacher框架结构图}
    \label{fig:placeholder}
\end{figure}

具体而言，如图3.1所示，学生模型（Student）和教师模型（Teacher）接收同一个输入 $x$，并分别加上不同的噪声 $\eta$ 和 $\eta'$。随后，学生模型的 Softmax 输出将用于计算两部分损失：(1) 与真实标签（One-hot label）计算交叉熵损失；(2) 与教师模型的输出计算一致性损失 $J(\theta)$。总损失根据预设权重加权后，用于更新学生模型。

在学生模型权重更新后，教师模型的权重通过指数移动平均（EMA）策略进行更新。权重更新规则可表示为：
\begin{equation}
    \theta'_{t} = \alpha \theta'_{t-1} + (1 - \alpha) \theta_{t}
\end{equation}

其中，$\alpha$ 为 EMA 系数，通常取接近 1 的值。这种系数的设定使得教师模型能够提供更加平滑的输出，从而为学生模型提供高质量的伪标签信息。此外，Mean Teacher 框架还引入了一致性损失函数，用于约束学生模型与教师模型在输入数据经过扰动后的输出一致性。设学生模型输出为 $y$，教师模型输出为 $y'$，则一致性损失 $L_{\text{consistency}}$ 可定义为两者的均方误差：
\begin{equation}
    L_{\text{consistency}} = \| y - y' \|^2
\end{equation}

这进一步增强了对未标记数据的利用率。从理论角度来讲，该方法的有效性可归因于其对模型输出的平滑性假设以及熵最小化原则的结合。

\subsection{Mean Teacher的优势与局限}
Mean Teacher框架在半监督领域与语义分割任务中展现出诸多优势，同时也存在一定的局限性。

从优势方面来看，首先，该框架通过引入教师模型与学生模型的协同训练机制，能够有效利用未标记数据，从而在标注数据有限的前提下显著提升模型性能。其次，由于教师模型的权重更新采用指数移动平均的方式，这一点能一定程度上缓解某些特殊样本对模型产生较大影响，使其输出结果具有较高的稳定性和可靠性。此外，Mean Teacher框架通过加入一致性损失函数，增强了模型对数据扰动的鲁棒性，从而在一定程度上缓解了过拟合问题。

然而，尽管Mean Teacher框架在性能上表现优异，但其局限性也不容忽视。例如，该框架对计算资源的需求较高，尤其是在处理大规模数据集时，训练时间的增加可能成为制约其实际应用的主要瓶颈。此外，因为要同时训练教师和学生两套模型的参数，所以内存占用也是单个模型的两倍。

\section{UA-MT框架构建与不确定性引入}
\subsection{UA-MT框架设计思路}

针对传统 Mean Teacher 框架在处理复杂语义分割任务时的局限性，本文提出了 UA-MT 框架。该框架的核心设计理念在于将认知不确定性建模融入Mean Teacher 的一致性正则化机制中。

在标准的 Mean Teacher 方法中，学生模型被强制拟合教师模型的输出，但这种拟合往往是“盲目”的，即无论教师模型的预测是否准确，学生都被要求与其保持一致。为了解决这一问题，UA-MT 引入了不确定性感知机制。具体而言，框架利用蒙特卡洛 Dropout 估算教师模型预测的认知不确定性，并将该不确定性作为权重引入到一致性损失函数中。

这种设计使得学生模型在训练过程中能够“有选择地”学习：对于教师模型预测置信度高的区域，学生模型严格遵循其预测；而对于教师模型本身也难以确定的模糊区域，则降低其对一致性损失的贡献，从而避免错误引导。通过这种机制，UA-MT 能够更有效地利用无标签数据，提升模型在复杂场景下的鲁棒性和泛化能力。

\subsection{基于蒙特卡洛 Dropout 的认知不确定性建模}
在语义分割任务中，模型的不确定性往往源于对模糊边界或相似类别的认知匮乏。为了量化这种不确定性，本研究利用教师模型生成不确定性评估图。具体而言，我们在推理阶段保持 Dropout 层开启，对同一输入数据进行多次随机前向传播，采用蒙特卡洛采样策略来近似模型的后验分布。

\begin{figure}
    \centering
    \includegraphics[width=1.0\linewidth]{project_30_3d28165c.png}
    \caption{UA-MT框架结构图}
    \label{fig:placeholder}
\end{figure}

如图所示，该图为UA-MT框架图。具体而言，对于每一个输入体素，我们进行 $T$ 次前向传播。设第 $t$ 次传播中，模型预测该体素属于第 $c$ 类别的概率为 $p^c_t$。为了获得稳健的不确定性估计，我们首先计算该体素在 $T$ 次采样中的平均预测概率 $\mu_c$：
\begin{equation}
    \mu_c = \frac{1}{T} \sum_{t=1}^{T} p^c_t
\end{equation}

其中，$\mu_c$ 代表了模型对该体素属于类别 $c$ 的平均置信度。基于此平均概率分布，我们采用预测熵来量化不确定性。对于每一个体素，其不确定性得分 $u$ 定义为：
\begin{equation}
    u = -\sum_{c} \mu_c \log \mu_c
\end{equation}

该公式通过计算平均概率分布的熵来衡量模型的"犹豫"程度。当 $\mu_c$ 在各个类别间分布均匀时（即模型在多次预测中摇摆不定），熵值 $u$ 较高，表示高不确定性；反之，若模型始终倾向于某一类别，则 $u$ 较低。

通过对图像中每一个体素重复上述计算，最终可以构建出一个与输入数据空间维度一致的不确定性张量 $\mathcal{U}$。若输入数据的维度为 $H \times W \times D$，则不确定性张量表示为：
\begin{equation}
    \mathcal{U} \in \mathbb{R}^{H \times W \times D}
\end{equation}

\subsection{基于不确定性图的一致性损失}

在估计出不确定性图 $\mathcal{U}$ 的指导下，我们需要过滤掉高不确定性的预测，仅选择置信度较高的区域作为学生模型学习的目标。

特别地，对于半监督分割任务，我们将不确定性感知一致性损失 $\mathcal{L}_c$ 设定为教师模型和学生模型在确定性较高区域的体素均方误差（MSE）损失。该损失函数定义如下：
\begin{equation}
    \mathcal{L}_c(f', f) = \frac{\sum_{v} \mathbb{I}(u_v < H) \| f'_v - f_v \|^2}{\sum_{v} \mathbb{I}(u_v < H)}
\end{equation}

式中各符号含义如下：

$\mathbb{I}(\cdot)$ 为指示函数，当条件满足时值为 1，否则为 0，用来筛选高置信度像素点。

$f'_v$ 和 $f_v$ 分别是教师模型和学生模型在第 $v$ 个体素处的预测输出。

$u_v$ 是第 $v$ 个体素处的估计不确定性得分。

$H$ 是预设的阈值，用于确定哪些区域被筛选。

通过引入该阈值 $H$，损失函数仅在不确定性低于该阈值的体素上计算误差。随着训练过程中不确定性感知一致性损失的作用，学生模型和教师模型都可以专注于学习更可靠的知识，从而逐步降低模型的整体不确定性。

\section{实验设计}
\subsection{数据集与预处理}
本实验采用 2018 年左心房（LA）分割挑战赛公开数据集进行实验验证。该数据集包含80例 3D图像作为训练集，以及 20 例作为独立测试集。

为了验证框架在标注稀缺样本场景下的性能，我们将原始的 80 例训练数据划分为两部分：

有标签集（Labeled Set）： 仅选取 16 例样本作为有监督信号来源。

无标签集（Unlabeled Set）： 剩余的 64 例样本仅提供图像，不提供标签，用于半监督学习。

数据预处理阶段，我们对图像进行了归一化处理，并采用随机裁剪（Random Crop）策略，将输入块大小固定为 112×112×80 以适应网络输入。

\subsection{网络架构与训练策略}

本实验采用 3D V-Net 作为基础分割网络架构。为了引入认知不确定性，我们在 V-Net 的解码器部分引入了 Dropout 层，并在推理阶段保持开启状态以进行蒙特卡洛采样。

训练参数设置如下：

优化器： 随机梯度下降（SGD），动量设为 0.9，权重衰减为 0.0001。

学习率策略： 初始学习率设为 0.01，并采用阶梯式衰减策略。

Batch Size： 每 GPU 批次大小为 4。

迭代次数： 最大迭代次数设为 6000 次。

对比方法： 为了评估 UA-MT 的有效性，我构建了以下对比组：

Baseline (Supervised Only)： 仅使用 16 例有标签数据进行全监督训练。

Mean Teacher (MT)： 使用 16 例标签 + 64 例无标签数据，标准 Mean Teacher 框架。

UA-MT (Ours)： 使用 16 例标签 + 64 例无标签数据，引入不确定性感知机制的改进框架。

\section{对比与结果分析}
为了全面评估分割性能，我们采用了四个广泛使用的评价指标：

1. Dice 相似系数 (Dice Coefficient)： 衡量分割重叠度，值越大越好。

2. Jaccard 指数 (Jaccard Index)： 衡量交并比，值越大越好。

3. 95\% 豪斯多夫距离 (95\% HD)： 衡量表面最大不一致距离，值越小越好。

4. 平均表面距离 (ASD)： 衡量平均表面误差，值越小越好。

\begin{table}[htbp]
  \centering
  \caption{实验对比结果}
  \label{tab:exp_results_fullborder}
  % 在列定义中使用 |l|c|... 来添加竖线
  \begin{tabular}{|l|c|c|c|c|}
    \hline % 表格顶部的横线
    \textbf{实验设置 (Method)} & \textbf{Dice(\%)} & \textbf{Jaccard(\%)} & \textbf{ASD(vox)} & \textbf{95HD(vox)} \\
    \hline % 表头下方的横线
    16样本全监督 & 85.83 & 76.34 & 11.96 & 2.85 \\
    \hline % 数据行之间的横线
    80样本全监督 & 91.53 & 84.42 & 5.02 & 1.63 \\
    \hline % 数据行之间的横线
    16+64 半监督 MT & 87.25 & 77.74 & 12.73 & 3.31 \\
    \hline % 数据行之间的横线
    16+64 半监督 UA-MT & 88.18 & 79.09 & 12.32 & 3.27 \\
    \hline % 表格底部的横线
  \end{tabular}
\end{table}

结果分析：
半监督学习的有效性： 对比 Baseline (85.83\%) 与 Standard MT (87.24\%)，引入无标签数据后，Dice 系数提升了约 1.4\%。这证明了利用未标注数据能够有效补充有标签数据的不足，提升模型的泛化能力。

UA-MT 的显著优势： 将 Standard MT 与 UA-MT 进行对比，无论是从各个参数的角度来看，都可以清晰地看到引入不确定性机制带来的增益。

重叠度指标： UA-MT 的 Dice 系数达到 88.18\%，相比 Standard MT 提升了 0.94\%；Jaccard 指数提升了 1.34\%。这表明，通过引入认知不确定性对一致性损失进行加权，模型能够过滤掉教师网络在高不确定区域（如模糊边界）产生的噪声，从而获得更准确的分割掩膜。

边界精度指标： 在表面距离指标上，UA-MT 的 95\% HD 和 ASD 分别为 12.31 mm 和 3.26 mm。虽然提升幅度相对较小，但依然优于 Standard MT。这说明不确定性感知机制帮助模型在复杂边界处做出了更稳健的预测，避免了因盲目拟合错误伪标签而导致的边界抖动。

综上所述，UA-MT 方法通过利用未标注数据并结合不确定性建模，在仅使用 20\% 有标签数据（16/80）的情况下，取得了优于标准半监督方法的分割性能。

\chapter{基于Transformer的多尺度自适应表示学习模型}
\section{基础模型选择（U - Net）}
\subsection{U-Net结构解析}

\begin{figure}
    \centering
    \includegraphics[width=0.5\linewidth]{project_30_b4887c2d.png}
    \caption{U-Net网络架构图}
    \label{fig:placeholder}
\end{figure}

如图所示，U-Net作为一种经典的编码器-解码器结构网络，在医学图像分割领域展现了卓越的性能。

在结构上，其主要由编码器、解码器以及跳跃连接组成。其编码器部分通过多次卷积和池化操作逐步降低特征图的空间分辨率，同时增加通道数以提取更抽象的特征表示。解码器则通过上采样和卷积操作逐步恢复特征图的空间分辨率，并通过跳跃连接将编码器中对应层次的浅层特征与解码器中的深层特征进行融合。

这种对称设计与跳跃连接原理使得U-Net能够有效结合高低层次的信息，从而在分割任务中同时保持细节信息和对全局上下文的理解。此外，U-Net的简单性和高效性也为其在医学图像分割领域的广泛应用提供了重要支持，尤其是在处理复杂解剖结构时表现尤为突出。

\subsection{U-Net的优势与局限}
U-Net在处理医学图像时展现出诸多显著优势，其中最为突出的是其捕捉多尺度特征的能力。通过编码器中的多次下采样操作，U-Net能够提取不同分辨率下的特征信息，从而适应不同大小的目标区域。此外，跳跃连接机制使得网络能够在解码过程中充分利用浅层特征中的细节信息，这对于精确分割医学图像中的细小结构尤为重要。

然而，U-Net在建模长距离依赖关系方面存在一定局限性。由于其主要依赖卷积操作进行特征提取，而卷积操作的感受野有限，因此难以有效捕捉图像中的全局上下文信息。这种局限性在处理具有复杂空间分布的大型病灶或器官时尤为明显，可能导致分割结果中出现过分割或欠分割的现象。此外，对于体积大小由较大差别的训练数据以及测试数据，U-Net因其结构简单，参数较少，无法自适应地完成训练以及推理。

\section{Transformer引入与融合策略}
\subsection{Swin - Unet模型}
\begin{figure}
    \centering
    \includegraphics[width=0.5\linewidth]{project_30_33e9464b.png}
    \caption{Swin - Unet模型架构图}
    \label{fig:placeholder}
\end{figure}

如图所示，Swin-Unet模型是一种将Swin Transformer模块融入U-Net架构的创新尝试，旨在通过引入自注意力机制来增强模型对多尺度特征的捕捉能力。

具体而言，Swin-Unet在编码器部分使用了基于窗口的自注意力机制（Window-based Self-Attention），这种机制通过将特征图划分为多个不重叠的窗口并在每个窗口内计算自注意力，显著降低了计算复杂度，同时保留了长距离依赖关系的建模能力。此外，Swin-Unet还采用了移位窗口策略（Shifted Window），以增强不同窗口之间的信息交互，从而进一步优化多尺度特征的提取效果，其他各个部分如跳跃连接等均是参考U-net完成的。
整体框架的结构：

1. Patch Partition + Linear Embedding：

这部分主要是将原图片的输入转化为Transformer接受的格式。即通过Patch划分以及线性映射将W × H × A的图像转变为W/4 × H/4 × C。

2. Swin Transformer Block模块：

这部分为核心模块，在下文中详细介绍。

3. Patch Merging：

这部分主要是将得到的图像再次进行压缩。

即将大小为W × H × C的图像通过合并以及线性映射的方式，合并相邻2 × 2范围内的像素，而后通过线性映射将通道压缩为1/2。

具体格式转换为W × H × C -> W/2 × H/2 × 2C。

\begin{figure}
    \centering
    \includegraphics[width=0.5\linewidth]{project_30_23d87a57.png}
    \caption{Swin Transformer模块结构图}
    \label{fig:placeholder}
\end{figure}

Swin Transformer内部具体具体工作方式：

上半部分：基于窗口的多头自注意力模块

1. 层归一化：输入特征 $z^{l-1}$ 首先经过层归一化处理，这有助于稳定训练过程。

2. 基于窗口的多头自注意力：这是核心计算单元。特征图被划分为若干不重叠的局部窗口，自注意力计算仅在每个窗口内部独立进行。其计算公式为：
\begin{equation}
\hat{z}^{l} = \operatorname{W-MSA}(\operatorname{LN}(z^{l-1})) + z^{l-1}
\end{equation}

这里使用了残差连接，将 W-MSA 的输出与原始输入 $z^{l-1}$ 相加，有助于梯度反向传播。

3. 多层感知机：中间输出 $\hat{z}^{l}$ 再次经过层归一化，然后进入多层感知机。最终输出为：
\begin{equation}
z^{l} = \operatorname{MLP}(\operatorname{LN}(\hat{z}^{l})) + \hat{z}^{l}
\end{equation}

下半部分：基于窗口的多头自注意力模块

1. 层归一化：输入 $z^{l}$ 经过层归一化。

2. 基于移位窗口的多头自注意力：这是 Swin Transformer 的精髓。相比于上半部分的规则窗口，这里的窗口划分向右下方移动了半个窗口的大小。这种移位操作使得新的窗口跨越了上一轮的窗口边界，从而在连续模块之间建立起跨窗口的连接。其计算公式为：
\begin{equation}
z^{l+1} = \operatorname{SW-MSA}(\operatorname{LN}(z^{l})) + z^{l}
\end{equation}

3. 多层感知机：最后，经过残差连接后的特征再次进入多层感知机进行非线性变换，得到最终输出：
\begin{equation}
z^{l+2} = \operatorname{MLP}(\operatorname{LN}(z^{l+1})) + z^{l+1}
\end{equation}

通过这种基于窗口的多头自注意力 + 基于移位窗口的多头自注意力的交替结构，Swin Transformer 既保持了局部计算的高效性（线性复杂度），又通过移位窗口机制逐步捕获了全局的上下文信息。

\subsection{TransUnet模型}
\begin{figure}
    \centering
    \includegraphics[width=0.5\linewidth]{project_30_189afccc.png}
    \caption{TransUnet模型结构图}
    \label{fig:placeholder}
\end{figure}

如图所示，TransUNet 模型是医学图像分割领域中首个将 Transformer 引入 U-Net 架构的开创性工作。与 Swin-Unet 的全 Transformer 结构不同，TransUNet 采用了一种 “CNN-Transformer 混合架构”。

其核心设计理念在于结合两者的优势：利用 CNN（如 ResNet）强大的局部特征提取能力来弥补 Transformer 在底层细节上的不足，同时利用 Transformer 的全局自注意力机制来解决 CNN 感受野受限、无法捕捉长距离依赖的问题。

整个框架的结构：

1.CNN 编码器：

	这部分主要用于提取图像的局部特征。不同于直接将原图切块，TransUNet 首先使用一个标准的 CNN作为主干网络。
    
具体来说，它将 W×H×C的图像通过卷积层和下采样层，转化为一个尺寸较小但通道数很深的特征图。

2.  Linear Projection

	这部分是将 CNN 提取的特征图转化为 Transformer 能够处理的序列格式。这部分在Swin-Unet中已经提及，作用相近。
    
3.Transformer 编码器

	这是模型的核心，用于捕捉全局上下文信息。输入是上述生成的特征序列。
通过多层 Transformer Layer，让序列中的每一个 Patch 都能与序列中其他所有 Patch 进行交互，从而建立全图的长距离依赖关系。

Transformer Layer内部具体具体工作方式：

上半部分：自注意力模块

1. 层归一化：输入特征 $z^{l-1}$ 首先经过层归一化处理，这有助于稳定训练过程。

2. 基于窗口的多头自注意力：这是核心计算单元。特征图被划分为若干不重叠的局部窗口，自注意力计算仅在每个窗口内部独立进行。其计算公式为：
\begin{equation}
\hat{z}^{l} = \operatorname{MSA}(\operatorname{LN}(z^{l-1})) + z^{l-1}
\end{equation}

这里使用了残差连接，将 MSA 的输出与原始输入 $z^{l-1}$ 相加，有助于梯度反向传播。

下半部分：基于窗口的多头自注意力模块

1. 层归一化：输入 $z^{l}$ 经过层归一化。

2. 多层感知机：中间输出 $z^{l}$ 再次经过层归一化，然后进入多层感知机。最终输出为：
\begin{equation}
z^{l+1} = \operatorname{MLP}(\operatorname{LN}(z^{l})) + z^{l}
\end{equation}

通过这种 CNN 提取局部特征 + Transformer 建立全局依赖的混合结构，TransUNet 既保留了 CNN 对医学图像边缘和纹理的敏感性，又成功引入了 Transformer 强大的全局建模能力，有效解决了传统 U-Net 在处理复杂解剖结构时的局限性。但由于其是对全局所有单元进行注意力计算，所以参数数量多，计算复杂度大，耗时较长。
\subsection{DA - TransUnet改进模型}
\begin{figure}
    \centering
    \includegraphics[width=0.5\linewidth]{project_30_a2ce7766.png}
    \caption{DA - TransUnet模型结构图}
    \label{fig:placeholder}
\end{figure}

\begin{figure}
    \centering
    \includegraphics[width=0.5\linewidth]{project_30_71a43e39.png}
    \caption{DA-Block模块结构图}
    \label{fig:placeholder}
\end{figure}

如图所示，DA-TransUNet 模型是在 TransUNet 基础上的一种改进架构，旨在通过引入双重注意力机制来进一步提升医学图像分割的精度。

具体而言，DA-TransUNet 在原有的 CNN 编码器和 Transformer 编码器之间，以及跳跃连接路径上，嵌入了专门设计的 DA-Block。这种机制能够同时在空间维度和通道维度上对特征进行加权，从而让模型更专注于病灶区域的细节特征，抑制无关背景的干扰。相比于 Swin-Unet 的窗口注意力，DA-TransUNet 更强调显式的特征校准。

整个框架的结构与TransUnet相近,只有跳跃连接的部分进行了改进。

故仅对跳跃连接处进行说明：

1.输入与通道降维：

对应途中的Conv块，将通道数压缩为n/16。

2.分支一：位置注意力模块PAM：

\begin{figure}
    \centering
    \includegraphics[width=0.5\linewidth]{project_30_507752c1.png}
    \caption{PAM模块结构图}
    \label{fig:placeholder}
\end{figure}

如图，该模块旨在生成一个空间注意力图，表示任意两个单元之间的影响。

该模块首先会让输入分别经过三个卷积层，而后生成特征图B、C、D，它们均被重塑为C×N的矩阵，以便计算位置间的相关性。

接着进行空间注意力图计算，主要通过矩阵乘法以及Softmax归一化函数完成：
\begin{equation}
S_{ji} = \frac{\exp(B_i \cdot C_j)}{\sum_{i=1}^{N} \exp(B_i \cdot C_j)}
\end{equation}

其中 $S_{ji}$ 表示第 $i$ 个位置对第 $j$ 个位置的影响程度（即注意力权重）。

最后将注意力图的转置与特征图 $D$ 进行矩阵乘法，并将结果重塑为原形状。在结尾处加入残差连接：
\begin{equation}
E_j = \alpha \sum_{i=1}^{N} (S_{ji} D_i) + A_j
\end{equation}

3.分支二：通道注意力模块CAM：

\begin{figure}
    \centering
    \includegraphics[width=0.5\linewidth]{project_30_6e81fa77.png}
    \caption{CAM模块结构图}
    \label{fig:placeholder}
\end{figure}

如图，这是CAM，它在提取通道特征发面表现优异，与PAM不同，直接将原始特征C×H×W重塑为C×N，然后对A和其转置进行矩阵乘法，接下来，使用一个Softmax层以获得通道注意力图。

4.特征融合与输出：

将上方支路和下方支路进行加权求和后，最后通过一个Conv层将图像恢复到目标大小以及通道数。

\section{实验设计}
\subsection{数据集与预处理}
本实验采用 Synapse 多器官分割数据集进行验证。该数据集包含腹部 CT 扫描图像，涵盖了胃、肾脏、脾脏、肝脏等 8 个主要器官以及背景，共计 9 个类别。

针对该数据集，进行了以下预处理：

图像重采样与裁剪：为了适配 Transformer 架构的输入要求，要将所有原始切片统一调整为 224×224的分辨率。

数据增强：在训练阶段，采用了随机旋转、翻转等在线增强策略，以增加样本的多样性，提升模型的泛化能力。

标准化：对图像像素强度进行归一化处理，消除不同扫描设备带来的分布差异。

\subsection{训练策略}

实验基于 PyTorch 框架，对比了四种不同的网络架构：UNet、TransUNet、Swin-Unet以及 DA-TransUNet。

所有模型均在相同的超参数配置下进行训练，以确保对比的公平性，具体训练参数设置如下：

优化器：Adam，利用其自适应学习率特性加速收敛。

学习率策略：初始学习率设为 0.01。为了平衡不同批次大小对梯度的影响，实验中根据公式 base\_lr = base\_lr * (batch\_size / 24) 对学习率进行了调整，使其根据不同的batch\_size自适应调整学习率。、

Batch Size：受硬件性能影响，每批次处理图像数量设为 12。

迭代次数：最大训练迭代次数设为 30,000 次（max\_iterations=30000），最大训练轮数设为 200。

输入尺寸：网络输入图像尺寸固定为 224×224 。

\section{对比与结果分析}
为了全面评估分割性能，实验采用了四个广泛使用的评价指标：

1. Dice 相似系数 (Dice Coefficient)： 衡量分割重叠度，值越大越好。

2. 95\% 豪斯多夫距离 (95\% HD)： 衡量表面最大不一致距离，值越小越好。

\begin{table}[htbp]
  \centering
  \caption{不同模型架构在测试集上的性能对比}
  \label{tab:model_comparison}
  \begin{tabular}{|l|c|c|}
    \hline
    \textbf{模型架构} & \textbf{Dice(\%)} & \textbf{95HD(mm)} \\
    \hline
    UNet & 76.12 & 39.10 \\
    \hline
    TransUNet & 76.69 & 24.61 \\
    \hline
    Swin-Unet & 77.93 & 31.09 \\
    \hline
    DA-TransUNet & 80.67 & 26.06 \\
    \hline
  \end{tabular}
\end{table}

结果分析：

基线模型表现：UNet 作为纯卷积网络，取得了 0.7612 的 Dice 系数，但在边界距离（HD95=39.10）上表现一般，说明其在捕捉长距离依赖和精确边界定位上存在局限。

Transformer 的引入：TransUNet 虽然引入了全局注意力机制，但在本实验设置下，其 Dice 系数（0.7669）仅略高于 UNet，不过其 HD95（24.61）显著优于 UNet，表明 Transformer 结构有助于改善边缘分割的准确性。

纯 Transformer 架构：Swin-Unet 展现了较强的特征提取能力，其 Dice 系数达到 0.7793，优于前两者，证明了基于滑动窗口的自注意力机制在处理医学图像纹理细节上的优势。

而对比两种主要使用Transformer结构的不同模型架构，TransUNet在边缘分割的准确性上效果较好，而Swin-Unet在图像重合度方向更胜一筹。但由于TransUNet需要对更多模型参数进行训练，所以耗费的时间远高于Swin-Unet，故在实际应用中应根据实际需要进行考量利用。

DA-TransUNet 在 TransUNet 的基础上进一步引入了注意力模块，从实验结果看，其Dice精度更高，而在边缘分割准确度方面与TransUNet相近，表现出了较好的性能。

\chapter{总结与展望}
医学图像分割是计算机辅助诊断系统中的关键环节，其精度直接影响医生的诊断效率与治疗方案。针对医学图像标注成本高昂以及传统卷积神经网络在长距离依赖建模上的局限性，本文分别从半监督学习策略与网络架构改进两个维度展开了深入研究。本文首先引入认知不确定性机制以利用大量无标签数据，随后设计了基于双重注意力机制的Transformer混合架构以提升分割性能

\section{论文所完成的工作}

本文围绕医学图像分割中的关键问题，主要完成了以下工作：

1.对深度学习在医学图像分割中的应用进行了系统性综述：

详细阐述了从全监督学习到半监督学习的范式转变，分析了卷积神经网络（CNN）与视觉Transformer（ViT）的各自优势与局限。重点探讨了Mean Teacher框架及认知不确定性（Epistemic Uncertainty）在解决小样本标注问题上的理论依据，为后续研究奠定了理论基础。

2. 提出了基于认知不确定性的半监督分割方法（UA-MT）

针对医学图像标注稀缺的问题，构建了UA-MT框架。

机制引入：在3D V-Net的解码器中引入Dropout层，利用蒙特卡洛采样估算模型的认知不确定性。

策略优化：在Mean Teacher架构的基础上，设计了不确定性加权的一致性损失函数，降低了无标签数据中不可靠区域的干扰。

实验验证：在Synapse数据集上验证了该方法仅需少量有标签数据（如16例）配合大量无标签数据，即可达到接近全监督模型的性能。

3. 构建了基于Transformer的多尺度自适应表示学习模型（DA-TransUNet）

针对纯CNN感受野受限和纯Transformer局部细节丢失的问题，本文深入研究了Swin-Unet与TransUNet架构，并在此基础上提出了改进模型DA-TransUNet。

架构融合：采用CNN与Transformer混合编码策略，既保留了CNN提取局部特征的能力，又利用Transformer捕获全局上下文信息。

双重注意力机制：在跳跃连接或瓶颈层引入了双重注意力模块，增强了模型对多尺度特征的自适应表达能力。

性能提升：实验结果表明，DA-TransUNet在Synapse多器官分割任务中，Mean Dice系数达到了 80.67\%，HD95指标显著降低，优于UNet、Swin-Unet及TransUNet等主流模型。

4. 完成了全面的对比实验与结果分析

在Synapse数据集上实施了严格的预处理与训练策略。通过定量指标（Mean Dice, Mean HD95）与定性可视化结果，充分证明了本文所提方法在处理复杂器官形态和边界模糊问题上的有效性。

\section{问题与不足}
尽管本文提出的方法在一定程度上提升了医学图像分割的精度，但受限于实验环境与算法复杂度，仍存在以下不足之处：

1. 计算资源消耗较大

	引入Transformer架构（尤其是Swin-Unet和DA-TransUNet）显著增加了模型的参数量和计算复杂度。在训练阶段，相比于传统的UNet，显存占用更高，推理时间更长，这在一定程度上限制了其在实时临床环境中的部署应用。
    
2. 超参数敏感性

	在UA-MT方法中，不确定性阈值的设定以及一致性损失的权重系数对最终结果影响较大，目前主要依赖经验进行调整，缺乏自适应的调节机制。在DA-TransUNet中，注意力模块的引入也增加了超参数搜索的空间。
    
3. 小器官分割仍存在误差

	从实验结果（如HD95指标）可以看出，虽然平均指标有所提升，但对于体积较小或形态变化剧烈的器官（如胆囊、胰腺），模型仍存在一定的漏检或边界分割不准现象。这可能与数据集的类别不平衡以及CT图像的切片厚度有关。
    
4. 数据模态单一

	本文仅在腹部CT数据集以及左心房CA数据集上进行了验证，未涉及MRI或其他模态的医学图像。不同模态的图像特性差异较大，本文模型的跨模态泛化能力尚需进一步验证。

\section{前景与展望}
基于本文的研究工作与存在的不足，未来的研究可以从以下几个方向展开：

1. 模型轻量化与加速

	为了推动深度学习模型在临床的落地，未来可研究基于知识蒸馏或模型剪枝的技术，压缩Transformer架构的参数量。探索轻量化注意力机制（如线性注意力）在医学分割中的应用，以在保证精度的同时大幅提升推理速度。
    
2.引入自监督预训练

	目前的实验主要依赖随机初始化或ImageNet预训练权重。未来可以探索在大规模无标签医学数据上进行自监督预训练（如MAE, SimMIM），学习更具鲁棒性的特征表示，从而进一步提升半监督学习的效果。
    
3.多模态融合与三维连续性建模

	目前的DA-TransUNet主要基于2D切片处理，忽略了层间的空间连续性。未来可尝试将模型扩展至3D版本，并探索CT与MRI等多模态数据的融合分割，利用不同模态的互补信息提升分割精度，特别是针对小器官的分割性能。
    
4.交互式分割与临床辅助

	结合本文的不确定性估计能力，开发人机交互式分割系统。当模型对某些区域预测不确定性较高时，自动提示医生进行点击或涂鸦修正，从而实现“人在回路”的高效精准分割，更好地服务于临床诊断。



\end{document}
